\chapter{Numerical validation in Python and Rust (Tier X)}\label{ch:numerics}

Chapters~1 and~2 stated upstream's theorem (AI-generated, not peer reviewed; kernel-checked in Lean, conditional on
the fidelity of the definitions) that the triangular lattice $\Lat$ has the least lower energy among all
locally finite configurations of density one, for every admissible potential $g$. This chapter does something much
weaker, and the weakness is the subject of its first section. We put a few hundred particles in a periodic box, let
them interact through a screened Coulomb (Yukawa) law, let the energy run downhill, and look at where it ends. The
code is the Rust crate \texttt{qf-yukawa-crystal} in rusty-SUNDIALS~\cite{rusty} (PR~\#72), and a short Python
companion written for this book. Every number in the chapter is \tierX{} unless a different tier is shown.

All Rust code below is quoted verbatim from the files
\texttt{crates/\allowbreak qf-yukawa-crystal/\allowbreak src/\allowbreak \{lib.rs,\allowbreak jellium.rs,\allowbreak main.rs\}} on the branch
\texttt{physics/\allowbreak jellium-\allowbreak wigner-\allowbreak 2026-10-11} of rusty-SUNDIALS, pinned to commit
\texttt{b7e4380} of that branch, with the line numbers given in each listing title. Every listing and every line
number in this chapter was re-checked against \texttt{git show b7e4380:<path>}.
Lines left out are marked \texttt{// [...]} (an insertion of ours, not part of the file).

\section{What numerics can and cannot do}\label{sec:num-can}

\subsection{The gap between the theorem and the computer}
The theorem of Chapter~2 is about objects that a computer cannot hold.
\begin{itemize}
\item \emph{Infinite configurations.} $C$ is an arbitrary locally finite subset of $\R^2$, and its lower energy is a
$\liminf$ over disks of radius $R\to\infty$. A simulation holds finitely many particles, here $N=288$.
\item \emph{All competitors.} The theorem quantifies over every $C$. A simulation visits the few end states that a
particular algorithm reaches from a few hundred starting points.
\item \emph{Exact potentials.} The theorem uses $g$ exactly. The code truncates the potential at a cutoff $r_c$ and
modifies it near $r_c$ (Section~\ref{sec:num-model}). The truncated potential is \emph{not} completely monotone, so the
theorem does not even apply to it.
\item \emph{Exact arithmetic.} The code uses IEEE double precision, with relative rounding $\approx 1.1\times10^{-16}$ per
operation and sums of tens of thousands of terms.
\end{itemize}
A periodic box is itself a competitor of a special kind. A configuration with $N$ particles in a periodic rectangle
of area $N$ is, after periodic extension, a locally finite configuration of density one in the plane. Its energy per
particle in the box (with the minimum-image rule and $r_c$ less than half the box) equals one half of the lower energy
$E_g$ of Chapter~1 for the periodic extension (Chapter~1 counts ordered pairs; this chapter, like the physics
literature, counts each pair once), for the truncated potential. So the box computes honest energies of honest competitors, but only of
\emph{periodic} ones, with one fixed period cell, and only for the truncated potential.

\begin{warnbox}[: Tier X never supports a claim]
In the Elenchus scheme~\cite{elenchus} a floating-point experiment is \tierX: it may steer a search, suggest a
conjecture or expose a bug, but it is never evidence for a theorem. ``No relaxation went below the lattice'' is
consistent with Chapter~2. It does not support it: the theorem is already kernel-checked (conditional on the fidelity
of its definitions), and a few hundred periodic states say nothing about the infinitely many non-periodic ones.
\end{warnbox}

What, then, is the point? There are four honest uses.
\begin{enumerate}
\item \emph{Testing code against a theorem.} If a relaxation ever ended below the lattice energy, the first suspect
would be the code, the second the truncation, and only a distant third the mathematics. A theorem is a strong test
oracle for numerical software.
\item \emph{Exploring the landscape.} The theorem identifies the minimum \emph{value}, not the local minima, their
depths or basins. Numerics can explore those.
\item \emph{Checking statement adequacy.} If the Lean definitions of Chapter~1 described the wrong object, the energies
computed from the same formulas would disagree with physical intuition or with the literature (for example the
Madelung constants of Section~\ref{sec:num-jellium}).
\item \emph{Connecting to physics} (Chapter~4): what finite systems at finite temperature actually do.
\end{enumerate}

\subsection{Preregistration and commit order}
An experiment whose success criterion is chosen after the results are seen can always be declared a success. The study
in this chapter was therefore \emph{preregistered}: its predictions, tolerances, potentials, box, initial conditions and
analysis were written into \texttt{results/preregistration.json} and committed (rusty-SUNDIALS commit
\texttt{636c514}) \emph{before} the simulation code (\texttt{968b0f2}) and before the results (\texttt{1f46e70}). The
order of commits is the evidence: anyone can check with \texttt{git log} that the criteria predate the code. The three
predictions read, verbatim:
\begin{quote}\small
\textbf{P1} \texttt{no\_undercut}: ``Over all runs (T = 0 quenches and annealed runs, all initial conditions), no
relaxed configuration has e < e\_tri - 1e-9 * |e\_tri| (kappa = 2). A violation would contradict the theorem (or reveal
a truncation artefact, to be checked at kappa = 1 and with larger r\_c).''\\[0.3em]
\textbf{P2} \texttt{annealed\_reaches\_lattice}: ``At least one annealed run per initial-condition family reaches the
perfect lattice: zero defects and |e - e\_tri| < 1e-8 |e\_tri|. Expected but not guaranteed (annealing is heuristic);
failure is reported, not hidden.''\\[0.3em]
\textbf{P3} \texttt{quench\_glassy}: ``T = 0 quenches from disordered starts mostly end in defected local minima (e >
e\_tri, defects > 0). This is a property of gradient flow, not evidence against the theorem.''
\end{quote}
P1 is the only prediction connected to the theorem. P2 is a prediction about an algorithm (annealing), and P3 about
another (gradient descent). A failed P2 is reported as failed (Section~\ref{sec:num-results}); it says nothing about the
theorem.

The P1 wording contains a hedge worth noticing: the potential is truncated, so a genuine undercut \emph{could} be a
truncation artefact. In the preregistration this is not an escape hatch but a test plan, fixed in advance: an undercut
would have to be re-examined at another $\kappa$ and a larger $r_c$.

\subsection{Controls must be shown to fire}
A check that cannot fail carries no information. Before any production number is reported, every observable is run on
inputs for which the right answer is known, in both directions.
\begin{itemize}
\item A \emph{positive control} is an input on which the check must pass: the perfect lattice must stay put, a slightly
perturbed lattice must relax back to it.
\item A \emph{negative control} is an input on which the check must fire: the square lattice must be detected as
defected and as a saddle, and a box that cannot hold the triangular lattice must produce defected states above
$e_{\rm tri}$.
\end{itemize}
If the defect counter reported zero defects for the square lattice, every ``zero defects'' in the results would be
worthless. The same goes for the Hessian test and for the energy comparison. This is the numerical twin of the negative
controls of the Lean work (\lean{unscaled_false}, \lean{riesz_neg_two_not_admissible}): an auditor that accepts
everything is not an auditor.

\section{The model}\label{sec:num-model}

\subsection{The Yukawa potential and complete monotonicity}
Two particles at distance $r$ interact through
\begin{equation}\label{eq:yukawa}
V(r)=\frac{e^{-\kappa r}}{r},\qquad \kappa>0 .
\end{equation}
In the language of Chapter~1, the potential is $g(t)=V(\sqrt t)=e^{-\kappa\sqrt t}\,t^{-1/2}$ as a function of the
squared distance $t=r^2$.

\begin{proposition}\label{prop:yukawa-cm}
For every $\kappa\ge0$, $g(t)=e^{-\kappa\sqrt t}\,t^{-1/2}$ is completely monotone on $(0,\infty)$.
\end{proposition}
\begin{humanproof}
\tierC{} (a standard argument, not formalised here). The function $t^{-1/2}$ is completely monotone:
$(-1)^r\frac{d^r}{dt^r}t^{-1/2}=\frac12\cdot\frac32\cdots\frac{2r-1}2\,t^{-1/2-r}\ge0$; this is the $s=1$ case of
\lean{riesz_admissible}. The function $h(t)=\kappa\sqrt t$ has a completely monotone derivative, and the composition
$e^{-h}$ of the completely monotone $u\mapsto e^{-u}$ with such an $h$ is completely monotone (differentiate
$e^{-h}$ repeatedly: every term is a product of $e^{-h}$ and derivatives of $-h$, each with the right sign). The product
of two completely monotone functions is completely monotone by the Leibniz rule, since every term of
$(-1)^r(fg)^{(r)}=\sum_k\binom rk\,(-1)^kf^{(k)}\,(-1)^{r-k}g^{(r-k)}$ is non-negative.
\end{humanproof}

So the theorem predicts that no density-one configuration has lower energy per particle than $\Lat$ for the exact
Yukawa potential. The values $\kappa=2$ (main) and $\kappa=1$ (check) are used throughout.

\subsection{Shifted-force cutoff}
The sum of $V$ over a lattice converges quickly, but a simulation needs a finite neighbour list. We cut at $r_c=8$
and \emph{shift the force}:
\begin{equation}\label{eq:vsf}
V_{\rm sf}(r)=\begin{cases}V(r)-V(r_c)-(r-r_c)\,V'(r_c), & r<r_c,\\ 0, & r\ge r_c.\end{cases}
\end{equation}
Then $V_{\rm sf}(r_c)=0$ and $V_{\rm sf}'(r_c^-)=V'(r_c)-V'(r_c)=0$: energy and force are continuous at $r_c$. The second
derivative is \emph{not} shifted, $V_{\rm sf}''=V''$ for $r<r_c$, so the Hessian jumps by $V''(r_c)$ when a pair
crosses $r_c$. The derivatives are
\begin{equation}\label{eq:vprime}
V'(r)=-\frac{e^{-\kappa r}(\kappa r+1)}{r^2},\qquad V''(r)=\frac{e^{-\kappa r}(\kappa^2r^2+2\kappa r+2)}{r^3}.
\end{equation}
In the code (\texttt{lib.rs}, lines 42--63):
\begin{lstlisting}[language={},title={lib.rs, lines 42--63}]
    fn v(&self, r: f64) -> f64 {
        (-self.kappa * r).exp() / r
    }
    fn dv(&self, r: f64) -> f64 {
        -(-self.kappa * r).exp() * (self.kappa * r + 1.0) / (r * r)
    }
    fn d2v(&self, r: f64) -> f64 {
        let k = self.kappa;
        (-k * r).exp() * (k * k * r * r + 2.0 * k * r + 2.0) / (r * r * r)
    }
    /// Shifted-force pair potential and its first two derivatives (zero for `r >= rc`).
    pub fn pair(&self, r: f64) -> (f64, f64, f64) {
        if r >= self.rc {
            return (0.0, 0.0, 0.0);
        }
        let (vc, dvc) = (self.v(self.rc), self.dv(self.rc));
        (
            self.v(r) - vc - (r - self.rc) * dvc,
            self.dv(r) - dvc,
            self.d2v(r),
        )
    }
\end{lstlisting}
The \texttt{refs} command records $V(r_c)=e^{-\kappa r_c}/r_c$ as \texttt{truncation\_v\_rc}:
$e^{-16}/8\approx1.41\times10^{-8}$ at $\kappa=2$ and $e^{-8}/8\approx4.19\times10^{-5}$ at
$\kappa=1$.\footnote{\texttt{python -I -c "import math; print(math.exp(-16)/8, math.exp(-8)/8)"}, output
\texttt{1.406689683990739e-08 4.193282848781398e-05}.} This is the pointwise value of $V$ at the cut, \emph{not} the
change in the energy. The linear shift $-V(r_c)-(r-r_c)V'(r_c)$ acts on every pair inside $r_c$, and there are many such
pairs (about $\pi r_c^2\approx200$ neighbours per particle), so the energy per particle moves much more. For the
triangular lattice, a scratch sum of the \emph{unshifted} $V$ out to $r=60$ gives $0.3935428489$ at $\kappa=2$ and
$1.6253213396$ at $\kappa=1$, against the shifted-force values $0.3935333668$ and $1.6076092172$ of
Table~\ref{tab:num-refs}: differences $9.5\times10^{-6}$ and $1.77\times10^{-2}$.\footnote{\label{fn:unshifted}Scratch
double loop over $|i|,|j|\le200$ (not shipped), $\frac12\sum_{0<|v|<60}e^{-\kappa|v|}/|v|$. Output (triangular, square):
$\kappa=1$: \texttt{1.6253213396062602}, \texttt{1.6348554953171857}; $\kappa=2$: \texttt{0.3935428489323898},
\texttt{0.40097424876505755}.} At $\kappa=1$ the shift is larger than $e_{\rm sq}-e_{\rm tri}$ itself. It is almost the
same for the two lattices, though: the unshifted square-minus-triangular differences are $9.53\times10^{-3}$ ($\kappa=1$)
and $7.43\times10^{-3}$ ($\kappa=2$), close to the shifted ones. Section~\ref{sec:num-python} shows the same effect when
$r_c$ goes from 6 to 8 at $\kappa=2$: the shifted-force lattice energy moves by $3.0\times10^{-4}$, while $V(6)$ is about
$10^{-6}$. What saves the study is that the lattice reference is computed with the same $V_{\rm sf}$: every energy
comparison in this chapter is between equally truncated energies.

\begin{remark}
$V_{\rm sf}$ is not completely monotone (it vanishes on $[r_c,\infty)$, is positive just below, and its second
derivative jumps), so P1 is a prediction for a potential the theorem does not cover. The modification is not small at
$\kappa=1$, only nearly constant per particle, which suggests (\tierC) that it preserves the ordering of
configurations.
\end{remark}

\subsection{The periodic box at density one}
At density one the triangular lattice has spacing $a$ with $a^2\sqrt3/2=1$, i.e.
\begin{equation}
a=\sqrt{2/\sqrt3}\approx1.074570 .
\end{equation}
The box is $L_x\times L_y=16a\times9\sqrt3\,a\approx17.1931\times16.7509$.\footnote{Evaluated with
\texttt{python -I -c "import math; a=math.sqrt(2/math.sqrt(3)); print(a, 16*a, 9*math.sqrt(3)*a)"}, output
\texttt{1.074569931823542 17.193118909176672 16.750887463837792}.} It holds $16$ sites per row and $18$ rows (rows
are $a\sqrt3/2$ apart, and $18\cdot a\sqrt3/2=9\sqrt3\,a$), so $N=288$ and the area is
$16\cdot9\sqrt3\,a^2=288$: density one, and the triangular lattice fits periodically. Such a box is called
\emph{commensurate} with $\Lat$. A square box of side $\sqrt{289}=17$ holds the square lattice ($N=289$) but not the
triangular one; it is used as a negative control.

Displacements use the minimum-image convention: \texttt{delta} (\texttt{lib.rs}, lines 65--72) replaces each coordinate
difference $d_x$ by \texttt{dx -= lx * (dx / lx).round()}, i.e.\ $d_x-L_x\,\mathrm{round}(d_x/L_x)\in[-L_x/2,L_x/2]$, and
likewise for $d_y$. The minimum image is the correct nearest copy as long as $r_c\le\min(L_x,L_y)/2$, which holds here ($8<8.375$). Under
this condition each pair contributes at most once, and the energy of the box is
\begin{equation}\label{eq:Ebox}
E(x)=\sum_{1\le i<j\le N}V_{\rm sf}\bigl(\lvert x_i-x_j\rvert_{\rm mi}\bigr),\qquad e(x)=E(x)/N .
\end{equation}
The minimum-image map jumps when a coordinate difference crosses $\pm L/2$, but there $|x_i-x_j|_{\rm mi}\ge L/2>r_c$
and $V_{\rm sf}=0$, so $E$ is continuously differentiable everywhere and smooth away from the pair distances $r_c$.

The energy per particle of the infinite lattice with basis $b_1,b_2$ is
\begin{equation}\label{eq:elat}
e_{\rm lat}=\frac12\sum_{v\in\Lambda\setminus\{0\}}V_{\rm sf}(|v|),
\end{equation}
which \texttt{lattice\_energy} (\texttt{lib.rs}, lines 224--239) computes by a direct double loop. For the commensurate
box the two must agree, and the unit test \texttt{periodic\_\allowbreak lattice\_\allowbreak energy\_\allowbreak equals\_\allowbreak direct\_\allowbreak lattice\_\allowbreak sum} checks this to
$10^{-12}$ relative; the study's references agree to about $10^{-14}$ (Table~\ref{tab:num-refs}).

\subsection{Force and Hessian}
Write $d=x_i-x_j$ (minimum image), $r=|d|$, $u=d/r$, and $\varphi(d)=V_{\rm sf}(|d|)$.

\begin{proposition}\label{prop:grad-hess}
For $0<r<r_c$,
\begin{equation}\label{eq:pairgrad}
\nabla_d\varphi=V_{\rm sf}'(r)\,u,\qquad
\nabla_d^2\varphi=K(d):=V''(r)\,uu^{\mathsf T}+\frac{V_{\rm sf}'(r)}{r}\bigl(I-uu^{\mathsf T}\bigr).
\end{equation}
Consequently the force on particle $i$ is $F_i=-\partial E/\partial x_i=-\sum_{j\ne i}V_{\rm sf}'(r_{ij})\,u_{ij}$, and
the Hessian $H=\nabla^2E\in\R^{2N\times2N}$ is the sum over pairs $i<j$ with $r_{ij}<r_c$ of the $2\times2$ blocks
\begin{equation}\label{eq:Hblocks}
H_{ii}\mathrel{+}=K,\qquad H_{jj}\mathrel{+}=K,\qquad H_{ij}\mathrel{-}=K,\qquad H_{ji}\mathrel{-}=K .
\end{equation}
\end{proposition}
\begin{humanproof}
$\partial r/\partial d_a=d_a/r=u_a$, so $\partial_a\varphi=V_{\rm sf}'(r)u_a$. Differentiating again,
$\partial_bu_a=(\delta_{ab}-u_au_b)/r$, hence
$\partial_b\partial_a\varphi=V_{\rm sf}''(r)u_au_b+V_{\rm sf}'(r)(\delta_{ab}-u_au_b)/r$, and $V_{\rm sf}''=V''$ below
$r_c$. Since $d=x_i-x_j$, $\partial d/\partial x_i=I$ and $\partial d/\partial x_j=-I$, which gives the signs in
\eqref{eq:Hblocks} and, for the gradient, $\partial E/\partial x_i=+\nabla_d\varphi$ and $\partial E/\partial x_j=-\nabla_d\varphi$.
\end{humanproof}

Three facts follow at once. (i) $\sum_iF_i=0$ (Newton's third law), so the gradient is orthogonal to rigid
translations. (ii) Each row of blocks of $H$ sums to zero, so the two translation vectors
$\tau_x=(1,0,1,0,\dots)$ and $\tau_y=(0,1,0,1,\dots)$ lie in $\ker H$. (iii) The radial part $V''uu^{\mathsf T}$ is
positive (since $V''>0$), while the transverse part $V_{\rm sf}'/r$ is negative (repulsion), which is what allows
saddles.

The force in the code, with the external (disk) terms trimmed:
\begin{lstlisting}[language={},title={lib.rs, lines 105--121 (external-potential part trimmed)}]
    /// Force `F = -grad E` written into `f`.
    pub fn force(&self, x: &[f64], f: &mut [f64]) {
        let n = x.len() / 2;
        f.iter_mut().for_each(|v| *v = 0.0);
        for i in 0..n {
            for j in (i + 1)..n {
                let (dx, dy) = self.delta(x[2 * i], x[2 * i + 1], x[2 * j], x[2 * j + 1]);
                let r = (dx * dx + dy * dy).sqrt();
                let (_, dv, _) = self.pair(r);
                if dv != 0.0 {
                    let (fx, fy) = (-dv * dx / r, -dv * dy / r);
                    f[2 * i] += fx;
                    f[2 * i + 1] += fy;
                    f[2 * j] -= fx;
                    f[2 * j + 1] -= fy;
                }
            }
            // [...] external potential U(r) for the disk geometries (lines 122-131)
\end{lstlisting}
and the Hessian, exactly \eqref{eq:Hblocks}:
\begin{lstlisting}[language={},title={lib.rs, lines 135--160 (external-potential part trimmed)}]
    /// Hessian of `E` (dense, row-major, `2n x 2n`).
    pub fn hessian(&self, x: &[f64]) -> Vec<f64> {
        let n = x.len() / 2;
        let m = 2 * n;
        let mut h = vec![0.0; m * m];
        let mut add = |a: usize, b: usize, v: f64| h[a * m + b] += v;
        for i in 0..n {
            for j in (i + 1)..n {
                let (dx, dy) = self.delta(x[2 * i], x[2 * i + 1], x[2 * j], x[2 * j + 1]);
                let r = (dx * dx + dy * dy).sqrt();
                let (_, dv, d2v) = self.pair(r);
                if r >= self.rc {
                    continue;
                }
                let u = [dx / r, dy / r];
                for a in 0..2 {
                    for b in 0..2 {
                        let id = if a == b { 1.0 } else { 0.0 };
                        let k = d2v * u[a] * u[b] + dv / r * (id - u[a] * u[b]);
                        add(2 * i + a, 2 * i + b, k);
                        add(2 * j + a, 2 * j + b, k);
                        add(2 * i + a, 2 * j + b, -k);
                        add(2 * j + a, 2 * i + b, -k);
                    }
                }
            }
            // [...] bowl / wall terms for the disk geometries (lines 161-183)
\end{lstlisting}
Formulas are checked against finite differences in the crate's unit tests:
\texttt{force\_\allowbreak matches\_\allowbreak finite\_\allowbreak difference\_\allowbreak of\_\allowbreak energy} (central difference $h=10^{-6}$, tolerance $10^{-6}$) and
\texttt{hessian\_\allowbreak matches\_\allowbreak finite\_\allowbreak difference\_\allowbreak of\_\allowbreak force} (tolerance $10^{-5}$), on a lattice perturbed by Gaussian noise of
standard deviation $0.05$. We ran the whole test suite once for this book (2 build jobs, 2 test threads):
\begin{lstlisting}[language={},title={cargo test -q --release -j 2 --offline -p qf-yukawa-crystal --manifest-path .../Cargo.toml -- --test-threads=2}]
running 10 tests
..........
test result: ok. 10 passed; 0 failed; 0 ignored; 0 measured; 0 filtered out; finished in 0.15s
[...] two further blocks "running 0 tests ... 0 passed; 0 failed"
\end{lstlisting}
The ten tests are the five of \texttt{lib.rs} (force and Hessian against finite differences, box energy against the
direct lattice sum, a Jacobi eigenvalue check on a matrix with known spectrum, square above triangular) and the five of
\texttt{jellium.rs} (Section~\ref{sec:num-jellium}). The line marked \texttt{[...]} is our summary of the omitted output (two further test targets with no tests, presumably
the binary and the doc-tests; \texttt{-q} hides their names).

\section{Gradient flow, stiffness and CVODE}\label{sec:num-flow}

\subsection{Relaxation as an ODE}
The relaxation is the \emph{gradient flow}
\begin{equation}\label{eq:gflow}
\dot x(t)=F(x(t))=-\nabla E(x(t)),\qquad x(0)=x_0\in\R^{2N}.
\end{equation}
Physically it is the overdamped motion of particles in a viscous fluid at zero temperature, which is how colloids move
(Chapter~4). Mathematically it has a Lyapunov function:
\begin{equation}
\frac{d}{dt}E(x(t))=\nabla E\cdot\dot x=-|\nabla E(x(t))|^2\le0 .
\end{equation}
The energy decreases strictly unless $x$ is a critical point. Since $E$ is bounded below on the box (it is a finite
sum of terms bounded below by $V_{\rm sf}\ge0$) and blows up when two particles collide, trajectories stay away from
collisions and approach the set of critical points. Generic starting points converge to local minima; a nondegenerate saddle attracts
only a measure-zero set of starting points (its stable manifold). A gradient flow ends where it ends: in whatever basin it started. This is
why P3 expects most quenches to end in defected minima.

Near a critical point $x^*$, the linearisation $\dot\delta=-H(x^*)\delta$ decays along each eigenvector of $H$ at rate
equal to its eigenvalue $\lambda_k$. The Jacobian of the right-hand side is
\begin{equation}
J(x)=\frac{\partial F}{\partial x}=-H(x).
\end{equation}

\subsection{Why stiff, and why BDF}
An ODE is \emph{stiff} when its Jacobian has eigenvalues of very different sizes, so that an explicit method must
take steps set by the fastest mode while the solution evolves on the slowest. For the perfect lattice at $\kappa=2$ we
computed the full spectrum of $H$ with the formula of Proposition~\ref{prop:grad-hess}: two zero modes, then the
smallest positive eigenvalue $2.355\times10^{-2}$ (twice), and the largest $5.700$, a ratio of about
$242$.\footnote{\label{fn:hess}A scratch check written for this chapter (not shipped with the book; it re-implements
\eqref{eq:Hblocks} in numpy and calls \texttt{numpy.linalg.eigvalsh}). Output:
\texttt{kappa=2.0 N=288 lowest=[-2.77382284e-15 2.77729229e-15 2.35534976e-02 2.35534976e-02]
largest=5.700435e+00 ratio\_max\_over\_third=2.4202e+02}. The lowest positive value agrees with the study's control
value $0.0236$.} At defected minima the study recorded positive eigenvalues as small as $5.4\times10^{-5}$
($\kappa=2$), so the slowest relaxation times are of order $1/(5.4\times10^{-5})\approx2\times10^4$ time units while the
fastest are below one. And far from equilibrium, in a random start where two particles sit at distance $0.5a$, the
local curvature $V''(0.5a)$ is larger still.

An explicit method such as forward Euler is stable for this linear problem only if the step satisfies
$h<2/\lambda_{\max}$; integrating to $t\sim10^5$ would need on the order of $10^5$ steps even where nothing happens.
Backward differentiation formulas (BDF) are implicit. They are A-stable for orders 1 and~2 and A($\alpha$)-stable for
orders 3 to~5, so in every case the whole negative real axis, where the eigenvalues $-\lambda_k$ of $J=-H$ lie, is
inside the stability region, and the step size is chosen by accuracy alone. The price
is a nonlinear solve per step,
\begin{equation}
x_{n+1}-h\beta_0F(x_{n+1})=\text{(known combination of past values)},
\end{equation}
solved by Newton's method with the matrix $I-h\beta_0J$. Here the Jacobian is available exactly: $J=-H$, from
Proposition~\ref{prop:grad-hess}. An exact Jacobian costs one Hessian evaluation ($O(N^2)$ pair work plus a dense
$2N\times2N$ matrix) but makes Newton converge in few iterations.

\subsection{CVODE in rusty-SUNDIALS}
rusty-SUNDIALS~\cite{rusty} is a Rust implementation of the SUNDIALS solver family. Its \texttt{cvode} crate provides
CVODE's variable-step BDF method of orders 1 to~5 in Nordsieck form (\texttt{constants.rs}, \texttt{nordsieck.rs}),
and forms the Newton matrix $M=I-\gamma J$ from a user-supplied dense Jacobian or, if none is given, from finite
differences (\texttt{solver.rs}). The relaxation driver, verbatim:
\begin{lstlisting}[language={},title={lib.rs, lines 252--300}]
/// Overdamped relaxation `dx/dt = F(x)` with CVODE (BDF, analytic Jacobian) until `max |F| < ftol` or `t > tmax`.
pub fn relax(sys: &System, x0: &[f64], ftol: f64, tmax: f64) -> Result<Relaxed, String> {
    let s1 = *sys;
    let s2 = *sys;
    let rhs = move |_t: f64, y: &[f64], ydot: &mut [f64]| -> Result<(), String> {
        s1.force(y, ydot);
        Ok(())
    };
    let jac = move |_t: f64, y: &[f64], j: &mut DenseMat| -> Result<(), String> {
        let h = s2.hessian(y);
        let m = y.len();
        for c in 0..m {
            for r in 0..m {
                j.cols[c][r] = -h[r * m + c];
            }
        }
        Ok(())
    };
    let mut solver = Cvode::builder(Method::Bdf)
        .rtol(1e-10)
        .atol(1e-12)
        .max_steps(2_000_000)
        .jacobian(jac)
        .build(rhs, 0.0, SerialVector::from_slice(x0))
        .map_err(|e| format!("CVODE setup failed: {e}"))?;
    let mut f = vec![0.0; x0.len()];
    let mut tout = 1.0;
    loop {
        let (t, y) = solver
            .solve(tout, Task::Normal)
            .map_err(|e| format!("CVODE failed at t = {tout}: {e}"))?;
        let y = y.to_vec();
        sys.force(&y, &mut f);
        let fmax = f.iter().fold(0.0_f64, |a, &b| a.max(b.abs()));
        if fmax < ftol || t >= tmax {
            let mut x = y;
            sys.wrap(&mut x);
            return Ok(Relaxed {
                x,
                t,
                max_force: fmax,
                steps: solver.num_steps(),
                rhs_evals: solver.num_rhs_evals(),
                converged: fmax < ftol,
            });
        }
        tout *= 2.0;
    }
}
\end{lstlisting}
Points to notice:
\begin{itemize}
\item The right-hand side is the force; the Jacobian callback writes $-H$ into CVODE's column-major dense matrix
(\texttt{j.cols[c][r]} is row $r$, column $c$).
\item Tolerances are tight ($10^{-10}$ relative, $10^{-12}$ absolute): we want the end point, not a pretty trajectory, but
a sloppy trajectory can cross into a different basin.
\item The solver is asked for output at $t=1,2,4,8,\dots$ (\texttt{tout *= 2.0}). The doubling matches the multiscale
spectrum: each output is a checkpoint, and the time horizon grows geometrically up to $t_{\max}=10^7$
(\texttt{TMAX} in \texttt{main.rs}).
\item The run stops when $\max_k|F_k|<10^{-9}$ (\texttt{FTOL} in \texttt{main.rs}). A run that hits $t_{\max}$ first is
returned with \texttt{converged: false}, and the summary counts it; it is not silently dropped.
\end{itemize}
Two absences, for the record: the rusty-SUNDIALS workspace has no KINSOL (a Newton solver for $F(x)=0$ directly), and
its \texttt{crates/ida} is a documented fixed-step prototype that the study does not use.

\subsection{What the stopping rule guarantees}
\begin{proposition}\label{prop:stop}
\tierC{} Suppose $E$ is exactly quadratic near a local minimum $x^*$, $E(x)=E(x^*)+\frac12\delta^{\mathsf T}H\delta$
with $\delta=x-x^*$, $H\succeq0$, $\ker H$ spanned by translations, and smallest positive eigenvalue $\lambda_*$. If the
gradient $g=\nabla E(x)=H\delta$ satisfies $\max_k|g_k|<\varepsilon$, then
\begin{equation}
0\le E(x)-E(x^*)\le\frac{|g|^2}{2\lambda_*}<\frac{2N\varepsilon^2}{2\lambda_*}.
\end{equation}
\end{proposition}
\begin{humanproof}
Translations do not change $E$, so we may take $\delta\perp\ker H$. Then $\delta=H^+g$ and
$E(x)-E(x^*)=\frac12g^{\mathsf T}H^+g\le|g|^2/(2\lambda_*)$, and $|g|^2\le2N\max_k|g_k|^2$.
\end{humanproof}
With $N=288$, $\varepsilon=10^{-9}$ and $\lambda_*\ge5.4\times10^{-5}$ this gives $E(x)-E(x^*)<5.3\times10^{-12}$, i.e.
less than $2\times10^{-14}$ per particle, far below the P1 tolerance $10^{-9}|e_{\rm tri}|\approx4\times10^{-10}$. The
argument is \tierC: real energies are not exactly quadratic and the bound uses the measured $\lambda_*$, but it shows that
the stopping rule is not the weak point.

\subsection{The Hessian spectrum as a minimum check}
A converged run is a critical point; it could still be a saddle. The second-order test settles it.
\begin{proposition}\label{prop:strictmin}
Let $E$ be $C^2$ near a critical point $x^*$, invariant under translations, with $H(x^*)\succeq0$ and $\ker H(x^*)$
equal to the translation plane $\mathrm{span}\{\tau_x,\tau_y\}$. Then $x^*$ is a strict local minimum modulo
translation: there is $\rho>0$ with $E(x)>E(x^*)$ for every $x$ with $0<\mathrm{dist}(x,x^*+\mathrm{span}\{\tau_x,\tau_y\})<\rho$.
\end{proposition}
\begin{humanproof}
Restrict $E$ to the affine slice $x^*+W$, $W=\mathrm{span}\{\tau_x,\tau_y\}^\perp$. There the Hessian is positive
definite, so $x^*$ is a strict local minimum on the slice (the usual second-order sufficient condition). Every point
near the orbit $x^*+\mathrm{span}\{\tau_x,\tau_y\}$ is a translate of a point of the slice near $x^*$, and $E$ is
translation invariant.
\end{humanproof}
On a torus, rotations are \emph{not} symmetries (they do not preserve the box), so for a periodic configuration we expect
exactly two zero modes. (For the disk with a harmonic bowl, translations are broken and a rotation is a symmetry, so
one expects one zero mode there; the disk runs were not Hessian-checked.)

The eigenvalues come from cyclic Jacobi rotations (\texttt{symmetric\_eigenvalues}, \texttt{lib.rs} lines 302--343),
which are slow ($O(m^3)$ per sweep, $m=576$) but simple, and accurate for small eigenvalues of symmetric matrices. The
crate checks it on the tridiagonal matrix with known eigenvalues $2-2\cos(k\pi/(m+1))$. The verdict on each run:
\begin{lstlisting}[language={},title={main.rs, lines 65--72}]
/// Lowest Hessian eigenvalues (after the two translation zero modes are included) for a periodic configuration.
fn hessian_summary(sys: &System, x: &[f64]) -> (f64, f64, f64, usize) {
    let m = x.len();
    let ev = symmetric_eigenvalues(&sys.hessian(x), m);
    let scale = ev[m - 1].abs();
    let neg = ev.iter().filter(|&&v| v < -1e-9 * scale).count();
    (ev[0], ev[1], ev[2], neg)
}
\end{lstlisting}
An eigenvalue counts as negative if it is below $-10^{-9}\lambda_{\max}$; the three lowest eigenvalues are recorded so
that one can see the two translation modes and the gap above them.

\subsection{Annealing, outside SUNDIALS}
A $T=0$ quench gets stuck. To give the lattice a chance, the study first \emph{anneals}: overdamped Langevin dynamics
\begin{equation}\label{eq:sde}
dX_t=F(X_t)\,dt+\sqrt{2T(t)}\,dW_t,
\end{equation}
with $W$ a standard Brownian motion in $\R^{2N}$ and a temperature lowered from $T_0=0.05$ to $T_1=10^{-5}$. At fixed
$T$ the stationary law is the Boltzmann distribution $\propto e^{-E/T}$, so slow cooling concentrates near low
minima; how slow is ``slow enough'' is exactly what the study tests (P2).

Equation~\eqref{eq:sde} is a stochastic differential equation, not an ODE. Brownian paths are nowhere differentiable,
so there is no smooth right-hand side for CVODE to interpolate, and CVODE's local error estimates, which compare
polynomial predictors with correctors, have no meaning. A stochastic scheme is needed. The code uses Euler--Maruyama,
$X_{k+1}=X_k+F(X_k)\,\Delta t+\sqrt{2T_k\Delta t}\,\xi_k$ with $\xi_k$ standard Gaussian, and a geometric schedule
$T_{k+1}=T_k(T_1/T_0)^{1/K}$:
\begin{lstlisting}[language={},title={lib.rs, lines 488--503}]
    let mut rng = Rng::new(seed);
    let mut x = x0.to_vec();
    let mut f = vec![0.0; x.len()];
    let ratio = (t1 / t0).powf(1.0 / steps as f64);
    let mut temp = t0;
    for _ in 0..steps {
        sys.force(&x, &mut f);
        let amp = (2.0 * temp * dt).sqrt();
        for k in 0..x.len() {
            let d = (f[k] * dt + amp * rng.gauss()).clamp(-cap, cap);
            x[k] += d;
        }
        sys.wrap(&mut x);
        temp *= ratio;
    }
    x
\end{lstlisting}
The study calls it with $K=300\,000$ steps, $\Delta t=0.01$ and \texttt{cap}$\,=0.05$
(\texttt{anneal(sys, x0, 0.05, 1e-5, 300\_000, 0.01, 0.05, seed)} in \texttt{main.rs}). Note the \texttt{clamp}: each
coordinate increment is capped at $0.05$ so that a close encounter cannot throw a particle across the box. With the cap
the scheme is a \emph{capped} Euler--Maruyama scheme, not a consistent discretisation of \eqref{eq:sde}; it is a
heuristic search, which is all P2 asks of it. Every annealed state is then quenched with CVODE as above, so the
reported end states are critical points of $E$, whatever the annealing did.

\section{Results of the preregistered study}\label{sec:num-results}

Every number in this section is \tierX{} and comes from the study's recorded results (rusty-SUNDIALS commit
\texttt{1f46e70}). Initial conditions are of three kinds: the perfect and noisy lattices (controls), uniform random
points with minimum separation $0.5a$, and points derived from the Hugging Face dataset
\texttt{colabfit/OC20\_IS2RES\_val\_id}~\cite{oc20} (CC-BY-4.0, shard \texttt{co/co\_10.parquet}): atomic structures
rotated at random in 3D, projected onto the plane and accepted with minimum separation $0.5a$. The dataset supplies
disorder only; nothing depends on its chemistry.

\subsection{References}
\begin{table}[h]\centering
\begin{tabular}{lccc}
\toprule
& $e_{\rm tri}$ & $e_{\rm sq}$ & $e_{\rm sq}-e_{\rm tri}$\\
\midrule
$\kappa=2$ & $0.3935333668$ & $0.4009647699$ & $7.43\times10^{-3}$\\
$\kappa=1$ & $1.6076092172$ & $1.6171455730$ & $9.54\times10^{-3}$\\
\bottomrule
\end{tabular}
\caption{Lattice energies per particle at density one, shifted-force Yukawa with $r_c=8$ (\tierX). The periodic-box
energy of the perfect lattice and the direct lattice sum agree to about $10^{-14}$. The Python companion of
Section~\ref{sec:num-python} reproduces $e_{\rm tri}(\kappa=2)$ independently.}\label{tab:num-refs}
\end{table}
The triangular lattice is below the square one at both screening lengths, as the theorem (and, for lattices alone, the
older results of Rankin, Cassels and Montgomery~\cite{rankin,cassels,montgomery} on Epstein zeta and theta functions)
predicts. Figure~\ref{fig:num-lattices} shows the references.

\begin{figure}[htbp]\centering
\includegraphics[width=0.8\textwidth]{lattices.pdf}
\caption{Reference energies per particle at density one (the vertical axes do not start at zero). Left and middle:
shifted-force Yukawa at the two values of $\kappa$, triangular below square. Right: the jellium extension of
Section~\ref{sec:num-jellium}. \tierX.}\label{fig:num-lattices}
\end{figure}

\subsection{Controls}
\begin{table}[h]\centering\small
\begin{tabular}{p{0.26\textwidth}p{0.30\textwidth}p{0.34\textwidth}}
\toprule
control & expected & observed\\
\midrule
perfect lattice (positive) & stays at $e_{\rm tri}$ & gap $5.2\times10^{-14}$\\
lattice $+$ noise $0.1a$ (positive) & relaxes back to $e_{\rm tri}$ & gap $2.5\times10^{-14}$\\
Hessian at the lattice (positive) & exactly 2 zero modes, then a gap & 2 zero modes, then $0.0236$\\
square lattice in $17\times17$ box (negative) & a saddle: negative eigenvalues & 72 negative eigenvalues ($\kappa=2$), 64 ($\kappa=1$)\\
incommensurate square box, $N=289$ (negative) & every relaxed state defected, above $e_{\rm tri}$ & all defected, minimum gap $9.6\times10^{-4}$ ($\kappa=2$)\\
\bottomrule
\end{tabular}
\caption{Positive and negative controls (\tierX). Each check is shown to pass where it must and to fire where it must
before any production number is read.}\label{tab:num-controls}
\end{table}
The negative controls matter most. The square lattice is a critical point of $E$ (by symmetry every particle feels zero
net force), so a code that only checked $\max|F|$ would accept it as an equilibrium. The Hessian test rejects it with
dozens of negative eigenvalues. The incommensurate box cannot hold the triangular lattice, and indeed no relaxed state
there reaches $e_{\rm tri}$ or is defect-free: the defect counter and the energy gap both fire.

\subsection{The three predictions}
\paragraph{P1 (no undercut): PASS.} At $\kappa=2$, 0 undercuts in 100 relaxations, 20 of them in the incommensurate
control box. At $\kappa=1$, 0 undercuts in 50 relaxations. (Each run of the main study is counted twice, as a $T=0$
quench and as an annealed-then-quenched state.) Figure~\ref{fig:num-gaps} shows every relaxation's gap.

\begin{figure}[htbp]\centering
\includegraphics[width=0.75\textwidth]{gaps.pdf}
\caption{Each relaxation's relative energy gap $(e-e_{\rm tri})/e_{\rm tri}$, one marker per run, on a logarithmic
axis. P1 concerns the sign of the gap: no gap is negative. Markers at the far left are runs that reached the lattice to
round-off. \tierX.}\label{fig:num-gaps}
\end{figure}

\paragraph{P2 (annealing reaches the lattice at least once per family).} At $\kappa=2$: PASS, narrowly. One of 20
OC20-derived and one of 20 uniform annealed runs reached the perfect lattice (zero defects and
$|e-e_{\rm tri}|<10^{-8}|e_{\rm tri}|$). At $\kappa=1$: \textbf{FAIL as preregistered}: 0 of 10 (OC20-derived) and 2 of 10
(uniform). The failure is a statement about the annealing schedule ($3\times10^5$ steps of capped Euler--Maruyama),
which is a weak instrument at $N=288$. It says nothing about the theorem in either direction, and it is reported, not
hidden, as the preregistration requires.

\paragraph{P3 (quenches are glassy): as expected.} At $\kappa=2$, 15 of 20 (OC20-derived) and 14 of 20 (uniform)
$T=0$ quenches end with defects.

\paragraph{Hessian coverage.} The final Hessian was computed for the 80 ($\kappa=2$) and 40 ($\kappa=1$) non-control
runs. None has a negative eigenvalue. The two translation modes satisfy $|\lambda|\le1.5\times10^{-15}$ ($\kappa=2$) and
$\le1.8\times10^{-15}$ ($\kappa=1$), and every other eigenvalue is at least $5.4\times10^{-5}$ ($\kappa=2$) and
$5.0\times10^{-5}$ ($\kappa=1$). By Proposition~\ref{prop:strictmin}, up to floating point, these end states are strict
local minima modulo translation. Ten orders of magnitude separate the zero modes from the smallest genuine eigenvalue,
so the classification is not borderline. The incommensurate-box and disk runs were not Hessian-checked.

Read together: many distinct local minima exist, every one found lies above the lattice, and annealing finds the
lattice only occasionally. The theorem provides none of this landscape information.

\subsection{Open versus closed disk (confounded)}
The Lean energy counts points in \emph{closed} disks (Chapter~1); the D4 work asks whether open disks give the same
answer. A tempting ``experiment'' is to confine $N=300$ particles ($\kappa=2$) in an open disk (harmonic bowl
$U=kr^2/2$, $k=0.02$) and in a closed disk (soft wall $U=\varepsilon(r-R)^4$ for $r>R$, $\varepsilon=100$), anneal and
quench, and compare. Results: interior defect fraction $0.086$ (open) versus $0.033$ (closed); edge defect fraction
$0.965$ versus $0.602$; interior mean $|\psi_6|$ $0.913$ versus $0.931$.

\begin{warnbox}[: a confounded comparison]
These numbers do not measure an open-versus-closed effect. The bowl produces a non-uniform density; the wall
compresses particles against the edge; and the two ``edge'' layers (within two spacings of the outermost particle)
contain different numbers of particles. The study describes finite-size edge effects in two different set-ups. It is
not a test of upstream's closed-disk counting, which is a question about measure-zero boundary terms in a limit $R\to\infty$
and is settled (or not) by proof, not by simulation.
\end{warnbox}

\begin{figure}[htbp]\centering
\includegraphics[width=0.7\textwidth]{disks.pdf}
\caption{Open (harmonic bowl) versus closed (soft wall) disk, $N=300$, $\kappa=2$: Voronoi cells coloured by
coordination, one run of each regenerated from its published seed. Defects concentrate at the edge in both cases.
\tierX, and confounded as stated in the text.}\label{fig:num-disks}
\end{figure}

The order parameters used throughout are those of \texttt{order} (\texttt{lib.rs}, lines 360--420): $d_0$ is the median
nearest-neighbour distance, two particles are neighbours if they are closer than $1.35\,d_0$, a particle is a defect if
it does not have exactly 6 neighbours, and $\psi_6(i)=\frac1{c_i}\sum_{j\sim i}e^{6\mathrm i\theta_{ij}}$ is the
bond-orientational order of particle $i$, with $\theta_{ij}$ the angle of the bond. For the perfect triangular lattice
$|\psi_6(i)|=1$; for the square lattice $\psi_6(i)=0$, since $\sum_{m=0}^3e^{6\mathrm i(m\pi/2)}=0$.

\section{The jellium extension: $s=1$ with Ewald summation}\label{sec:num-jellium}

\subsection{Why a background}
Here `Coulomb' means the three-dimensional $1/r$ law restricted to the plane (Riesz $s=1$); the planar Coulomb
(logarithmic, $s=0$) interaction is distinguished where it matters. For the Coulomb potential $1/r$ (Riesz $s=1$ in the plane) the lattice sum $\sum_{v\ne0}|v|^{-1}$ diverges: there are
about $2\pi r\,dr$ points at distance $r$, and $\int^\infty2\pi r\cdot r^{-1}\,dr=\infty$. Every competitor of
density one has infinite energy, and the theorem says nothing (\lean{latticeEnergy_riesz_lt_top} needs $s>2$, and
the inequality $\infty\le\infty$ is empty). Physics removes the divergence with a uniform neutralising background of
opposite charge: the \emph{one-component plasma}, or \emph{jellium}. Its energy per particle is finite, and for $s=1$ the
classical ground state is the two-dimensional Wigner crystal. Upstream's manuscript treats renormalized energies for
$0\le s<2$ in its Section~8, following Petrache and Serfaty~\cite{ps}, and concludes that the triangular lattice
minimizes them~\cite{oai} \tierL. This is \emph{not} formalised in Lean anywhere. The computation below is a numerical
first step, \tierX.

\subsection{The Ewald energy}
Let $N$ unit charges sit at $x_1,\dots,x_N$ in a periodic rectangle of area $\mathcal A=L_xL_y$, with background density
$n=N/\mathcal A$. Choose a splitting parameter $\alpha>0$ and write
\begin{equation}
\frac1r=\underbrace{\frac{\operatorname{erfc}(\alpha r)}{r}}_{\text{short range}}+\underbrace{\frac{\operatorname{erf}(\alpha r)}{r}}_{\text{smooth, long range}} .
\end{equation}
\begin{proposition}\label{prop:ewald}
With $S(k)=\sum_je^{\mathrm ik\cdot x_j}$ and $k$ ranging over the reciprocal lattice $\{(2\pi m/L_x,2\pi n'/L_y)\}$,
the energy of the neutral periodic system (particles plus background) is
\begin{equation}\label{eq:ewald}
E=\frac12{\sum_{i,j}}'\frac{\operatorname{erfc}(\alpha r_{ij})}{r_{ij}}
+\frac{\pi}{\mathcal A}\sum_{k\ne0}\frac{\operatorname{erfc}(k/2\alpha)}{k}\,|S(k)|^2
-\frac{N\alpha}{\sqrt\pi}-\frac{N^2\sqrt\pi}{\alpha\mathcal A},
\end{equation}
where $\sum'$ runs over all pairs and periodic images except $i=j$ in the same cell. $E$ does not depend on $\alpha$.
\end{proposition}
\begin{humanproof}
\tierC{} (the classical derivation; we give its steps.)
\emph{Long-range part.} The function $\operatorname{erf}(\alpha r)/r$ is smooth and its two-dimensional Fourier
transform is $\frac{2\pi}{k}\operatorname{erfc}\bigl(\frac{k}{2\alpha}\bigr)$. Summing it over all pairs and images
(including $i=j$) is a sum over the periodic charge density, which by the Poisson summation formula equals
$\frac1{2\mathcal A}\sum_k\frac{2\pi}k\operatorname{erfc}(\frac k{2\alpha})|S(k)|^2$. The $k\ne0$ terms give the
second term of \eqref{eq:ewald}. The $k=0$ term diverges, like $2\pi N^2/(2\mathcal A k)$, and is cancelled exactly by
the long-range interaction of the particles with the background and of the background with itself: a neutral system
has no $k=0$ Fourier component.
\emph{Self term.} The long-range sum included $i=j$ at distance $0$, which is not a real interaction. Its value
$\lim_{r\to0}\operatorname{erf}(\alpha r)/r=2\alpha/\sqrt\pi$, with the factor $\frac12$ and $N$ particles, gives
$-N\alpha/\sqrt\pi$.
\emph{Background in the short-range part.} The background also interacts through the short-range kernel. With
$\int_{\R^2}\frac{\operatorname{erfc}(\alpha r)}r\,d^2x=2\pi\int_0^\infty\operatorname{erfc}(\alpha r)\,dr=\frac{2\sqrt\pi}\alpha$,
the particle--background energy is $-Nn\frac{2\sqrt\pi}\alpha$ and the background--background energy is
$+\frac12Nn\frac{2\sqrt\pi}\alpha$; the sum is $-\frac{N^2\sqrt\pi}{\alpha\mathcal A}$.
\emph{Independence of $\alpha$.} Only the splitting of $1/r$ changes with $\alpha$, not the total.
\end{humanproof}

The code states the same formula in its header and implements it with a half-space $k$-sum (since $|S(-k)|=|S(k)|$,
each $k$ is taken once with weight~2):
\begin{lstlisting}[language={},title={jellium.rs, lines 4--9}]
//! Energy (with `A = Lx Ly`, `S(k) = sum_j exp(i k.x_j)`):
//! `E = 1/2 sum'_{i,j} erfc(a r_ij)/r_ij + (pi/A) sum_{k != 0} erfc(k/(2a))/k |S(k)|^2 - N a/sqrt(pi) - N^2 sqrt(pi)/(a A)`,
//! with the real-space sum over minimum images inside `rc` (requires `rc <= min(Lx, Ly)/2`).
//! The last term is what remains of the particle-background and background-background interactions after the
//! divergent `k = 0` parts cancel. This is the renormalized (jellium) energy of upstream's Section 8 for `s = 1`, up to
//! upstream's normalisation constant.
\end{lstlisting}
The last sentence of that comment is a \tierC{} identification (the normalisation constant is not computed here).
The constructor \texttt{Jellium::new} (lines 27--61) stores, for each half-space vector $k$ with $|k|<k_{\rm cut}$, the
coefficient \texttt{c = 2.0 * libm::erfc(k / (2.0 * alpha)) / k}; the factor 2 is the weight of the pair $\pm k$.
\begin{lstlisting}[language={},title={jellium.rs, lines 85--110}]
    /// Total energy.
    pub fn energy(&self, x: &[f64]) -> f64 {
        let n = x.len() / 2;
        let geo = self.geometry_system();
        let area = self.lx * self.ly;
        let mut real = 0.0;
        for i in 0..n {
            for j in (i + 1)..n {
                let (dx, dy) = geo.delta(x[2 * i], x[2 * i + 1], x[2 * j], x[2 * j + 1]);
                let r = (dx * dx + dy * dy).sqrt();
                if r < self.rc {
                    real += libm::erfc(self.alpha * r) / r;
                }
            }
        }
        let mut recip = 0.0;
        for &(kx, ky, c) in &self.kvecs {
            let (re, im) = self.structure(x, kx, ky);
            recip += c * (re * re + im * im);
        }
        recip *= std::f64::consts::PI / area;
        let nf = n as f64;
        real + recip
            - nf * self.alpha / std::f64::consts::PI.sqrt()
            - nf * nf * std::f64::consts::PI.sqrt() / (self.alpha * area)
    }
\end{lstlisting}
The real-space sum keeps only minimum images inside $r_c=8$ and the reciprocal sum only $|k|<k_{\rm cut}=7$, with
$\alpha=0.66$ (\texttt{tri\_jellium} in \texttt{main.rs}). Both truncations are tiny:
$\operatorname{erfc}(\alpha r_c)=\operatorname{erfc}(5.28)\approx8.2\times10^{-14}$ and
$\operatorname{erfc}(k_{\rm cut}/2\alpha)\approx6.4\times10^{-14}$.\footnote{Printed by the scratch check of
footnote~\ref{fn:hess}: \texttt{erfc(0.66*8) = 8.201410123512745e-14  erfc(7/(2*0.66)) = 6.400438678401495e-14}.}

The force follows by differentiating \eqref{eq:ewald}. The real-space pair force has magnitude
$-\frac{d}{dr}\frac{\operatorname{erfc}(\alpha r)}r=\frac{\operatorname{erfc}(\alpha r)}{r^2}+\frac{2\alpha}{\sqrt\pi}\frac{e^{-\alpha^2r^2}}{r}$.
For the reciprocal part, with $\phi_i=k\cdot x_i$,
$\partial_{x_i}|S(k)|^2=2\,\mathrm{Re}\bigl(\mathrm ik\,e^{\mathrm i\phi_i}\overline{S(k)}\bigr)=-2k\,\mathrm{Im}\bigl(e^{\mathrm i\phi_i}\overline{S(k)}\bigr)$
and $\mathrm{Im}(e^{\mathrm i\phi}\overline S)=\mathrm{Re}S\sin\phi-\mathrm{Im}S\cos\phi$. The code (lines 134--145)
implements exactly this, with the comment \texttt{// Im(conj(S) e\^{}\{i ph\}) = sre sin(ph) - sim cos(ph)} and the
prefactor \texttt{pref = 2.0 * std::f64::consts::PI / area}, and the unit test
\texttt{force\_\allowbreak matches\_\allowbreak finite\_\allowbreak difference} checks it against the energy.
Unlike the Yukawa case, the jellium relaxation does \emph{not} pass an analytic Jacobian: \texttt{Jellium::relax} builds
CVODE without \texttt{.jacobian(...)}, so CVODE forms its own difference-quotient Jacobian (the doc comment says
``finite-difference Jacobian''). The Hessian used for the spectrum checks is \texttt{hessian\_fd}: central finite
differences of the analytic force with step $10^{-5}$, symmetrised. Annealing uses the same capped Euler--Maruyama with
$10^5$ steps, because each force costs a $k$-sum.

\subsection{Controls}
The jellium extension was preregistered separately (rusty-SUNDIALS branch \texttt{physics/\allowbreak jellium-\allowbreak wigner-\allowbreak 2026-10-11};
preregistration \texttt{e706f85}, code \texttt{3bc55fa}, credit to Alves et al.\ \texttt{b7e4380}). Its controls:
\begin{itemize}
\item \emph{Madelung energies against the literature.} At density one ($r_s=1/\sqrt\pi$) the energy per particle is
$\eta\sqrt\pi$, with $\eta$ the Madelung constant in Hartree$\times r_s$. The crate uses the ten-digit values of Alves,
Bendazzoli, Evangelisti and Berger~\cite{alves}, which agree with Bonsall and Maradudin~\cite{bm}:
\begin{lstlisting}[language={},title={jellium.rs, lines 252--253}]
pub const ETA_TRIANGULAR: f64 = -1.106102587;
pub const ETA_SQUARE: f64 = -1.100244420;
\end{lstlisting}
The Ewald code gives $e_{\rm tri}=-1.960515789320$ against the literature value $-1.960515789825$ (relative difference
$2.6\times10^{-10}$), and $e_{\rm sq}=-1.950132460$ (relative difference $4.3\times10^{-10}$). This is the rounding level
of nine-decimal published constants. \tierL{} for the literature values, \tierX{} for ours.
\item \emph{Square saddle.} The square lattice has 64 unstable modes (negative eigenvalues of \texttt{hessian\_fd}),
consistent with the imaginary phonon branch reported by Bonsall and Maradudin~\cite{bm}.
\item \emph{$\alpha$-invariance.} Proposition~\ref{prop:ewald} says $E$ does not depend on $\alpha$. The test perturbs
the lattice by noise of standard deviation $0.2$ and compares three values of $\alpha$:
\begin{lstlisting}[language={},title={jellium.rs, lines 304--311}]
        // A 4x range of the splitting parameter, each with its k-space cutoff scaled so both truncations stay ~1e-13.
        let e: Vec<f64> = [0.66, 1.32, 2.64]
            .iter()
            .map(|&al| Jellium::new(lx, ly, al, 8.0, 10.6 * al).energy(&x))
            .collect();
        for v in &e[1..] {
            assert!(((v - e[0]) / e[0]).abs() < 1e-10, "{:?}", e);
        }
\end{lstlisting}
The energy changes by less than $10^{-10}$ relative over this $4\times$ range. A wrong self term or background term
would break this test, because those terms depend on $\alpha$ explicitly: it is a control that fires on the most
likely bugs.
\end{itemize}

\subsection{The main jellium relaxation study}
The preregistered prediction (J1) is the analogue of P1: no relaxed configuration below the triangular Madelung
energy. Status of the main study: \JelliumMainStatus.

\section{A from-scratch Python companion}\label{sec:num-python}

The Rust crate is the instrument of record; for learning, a script short enough to read in one sitting is better.
The script below, \texttt{code/ch3\_yukawa\_small.py} (78 lines, numpy and scipy only), builds the same shifted-force
Yukawa model in a small commensurate box, $n\times2n$ triangular sites, $N=2n^2$, area $N$, and minimises the energy
with L-BFGS from a perturbed lattice and from five random starts.

\begin{warnbox}[: L-BFGS is not CVODE]
L-BFGS is a quasi-Newton \emph{minimiser}: it takes line-search steps downhill using a low-rank curvature model. It
does not follow the gradient-flow trajectory \eqref{eq:gflow}, so from the same start it can end in a different local
minimum than CVODE. The script computes no Hessian, so it cannot distinguish a minimum from a saddle. Everything
it prints is \tierX.
\end{warnbox}

\lstinputlisting[language=Python,title={code/ch3\_yukawa\_small.py}]{code/ch3_yukawa_small.py}

One design point differs from the Rust code. At $n=6$ the box is $6.447\times11.167$, so $r_c=6$ (or $8$) exceeds
$L_x/2$, and the minimum image alone would miss pairs. The script therefore adds image shells, $M=\lceil r_c/L+\frac12\rceil$
copies in each direction beyond the minimum image, and sums over all images inside $r_c$, including a particle's own
images. The Rust code can use the minimum image only because its box is large enough ($8<8.375$). The gradient is
$\sum_{j,\text{images}}V_{\rm sf}'(r)\,d/r$; the factor $\frac12$ of the energy disappears because each pair appears twice
in the double sum (Exercise~\ref{ex:num-grad}).

We ran the script with its defaults ($n=6$, $\kappa=2$, $r_c=6$) on two cores (\texttt{taskset -c 2,3}); it took
26.8~s of wall-clock time. The exact output:
\begin{lstlisting}[language={},title={python -I code/ch3\_yukawa\_small.py}]
N=72 kappa=2.0 rc=6.0 box=6.447420x11.167258
e_tri box sum = 0.393231076156  direct lattice sum = 0.393231076156
perturbed lattice 0.1a   e-e_tri =  1.332e-15  max|grad| = 1.2e-07  iters =    62  defects = 0
random start seed 1      e-e_tri =  7.772e-16  max|grad| = 4.0e-08  iters =    95  defects = 0
random start seed 2      e-e_tri =  1.903e-03  max|grad| = 4.8e-08  iters =   165  defects = 4
random start seed 3      e-e_tri =  2.220e-16  max|grad| = 7.5e-08  iters =   112  defects = 0
random start seed 4      e-e_tri =  6.661e-16  max|grad| = 3.6e-08  iters =   146  defects = 0
random start seed 5      e-e_tri =  1.833e-03  max|grad| = 2.1e-07  iters =   222  defects = 4
\end{lstlisting}
And with the study's cutoff, \texttt{python -I code/ch3\_yukawa\_small.py 6 2 8}:
\begin{lstlisting}[language={},title={python -I code/ch3\_yukawa\_small.py 6 2 8}]
N=72 kappa=2.0 rc=8.0 box=6.447420x11.167258
e_tri box sum = 0.393533366826  direct lattice sum = 0.393533366826
perturbed lattice 0.1a   e-e_tri =  1.388e-15  max|grad| = 1.2e-07  iters =    62  defects = 0
random start seed 1      e-e_tri =  7.772e-16  max|grad| = 7.5e-08  iters =    94  defects = 0
random start seed 2      e-e_tri =  1.903e-03  max|grad| = 1.3e-07  iters =   160  defects = 4
random start seed 3      e-e_tri =  8.882e-16  max|grad| = 5.6e-08  iters =   114  defects = 0
random start seed 4      e-e_tri =  9.437e-16  max|grad| = 4.4e-08  iters =   145  defects = 0
random start seed 5      e-e_tri =  1.833e-03  max|grad| = 4.7e-08  iters =   221  defects = 4
\end{lstlisting}
How to read this output honestly:
\begin{itemize}
\item \emph{An independent reference check.} With $r_c=8$ the script's box energy and direct lattice sum,
$0.393533366826$, agree with the Rust reference $e_{\rm tri}(\kappa=2)=0.3935333668$ of Table~\ref{tab:num-refs}, from
an independent implementation in another language with a different image scheme and $N=72$ instead of $288$. (The
lattice energy per particle does not depend on $N$ for a commensurate box.) With $r_c=6$ the value is
$0.393231076156$: the cutoff changes the fourth significant digit at $\kappa=2$, which is why both codes compare
equally truncated energies.
\item \emph{The stopping rule was not met as asked.} The script requests \texttt{gtol=1e-10}, but the final
$\max|\nabla E|$ is between $4\times10^{-8}$ and $2\times10^{-7}$: L-BFGS stopped on its relative energy-decrease
criterion \texttt{ftol}. (A scratch rerun of two starts, the perturbed lattice and seed~2, printed scipy's message
\texttt{CONVERGENCE: RELATIVE REDUCTION OF F <= FACTR*EPSMCH} for both.) This is looser than the Rust rule $\max|F|<10^{-9}$. For the lattice runs the energy agrees with
$e_{\rm tri}$ to about $10^{-15}$ anyway, as Proposition~\ref{prop:stop} leads one to expect (the error is quadratic in
the gradient).
\item \emph{No undercut}, in twelve relaxations. Three of the five random starts (seeds 1, 3, 4) reached the
lattice (within $10^{-15}$, zero defects); two (seeds 2, 5) ended $1.8$--$1.9\times10^{-3}$ above it with 4 defects. Do
\emph{not} compare ``2 of 5 defected'' with P3 (15 of 20 quenches defected): system size ($72$ versus $288$), algorithm (L-BFGS versus gradient
flow) and cutoff all differ.
\item \emph{No designed negative control.} The two defected end states show that the defect counter and the energy gap
\emph{can} fire, but that is incidental, not designed. By the standard of Section~\ref{sec:num-can} the script is
incomplete until you add one; Exercises~\ref{ex:num-neg} and~\ref{ex:num-square} do that.
\end{itemize}

\section{How to reproduce the Rust study}\label{sec:num-repro}

All commands are run from the rusty-SUNDIALS checkout (branch \texttt{physics/\allowbreak jellium-\allowbreak wigner-\allowbreak 2026-10-11} for the jellium
commands). The binary is \texttt{yukawa\_crystal} (\texttt{[[bin]]} in the crate's \texttt{Cargo.toml}); every subcommand takes
the crate directory and writes JSON lines under \texttt{<crate\_dir>/results/} (\texttt{main.rs}, lines 1--3: the
``Experiment driver for the preregistered study \texttt{results/preregistration.json} (YUK-TRI-1)''). The number of worker threads is read from the environment:
\begin{lstlisting}[language={},title={main.rs, lines 260--265}]
    // Worker threads: YUKAWA_THREADS if set, else 2 (the machine is shared; the owner asked for 2 cores overnight).
    let threads = std::env::var("YUKAWA_THREADS")
        .ok()
        .and_then(|v| v.parse::<usize>().ok())
        .filter(|&t| t > 0)
        .unwrap_or(2);
\end{lstlisting}
The workflow:
\begin{lstlisting}[language={},title={Reproduction (one command per line; run from the repository root)}]
cargo test --release -p qf-yukawa-crystal
cargo build --release -p qf-yukawa-crystal

# initial conditions: OC20 shard co/co_10.parquet from colabfit/OC20_IS2RES_val_id
python crates/qf-yukawa-crystal/scripts/hf_initial_conditions.py co_10.parquet crates/qf-yukawa-crystal/data/initial

YUKAWA_THREADS=2 target/release/yukawa_crystal refs     crates/qf-yukawa-crystal
YUKAWA_THREADS=2 target/release/yukawa_crystal controls crates/qf-yukawa-crystal
YUKAWA_THREADS=2 target/release/yukawa_crystal main     crates/qf-yukawa-crystal 2 20
YUKAWA_THREADS=2 target/release/yukawa_crystal main     crates/qf-yukawa-crystal 1 10
YUKAWA_THREADS=2 target/release/yukawa_crystal disk     crates/qf-yukawa-crystal
YUKAWA_THREADS=2 target/release/yukawa_crystal snapshot crates/qf-yukawa-crystal
python crates/qf-yukawa-crystal/scripts/summarize.py crates/qf-yukawa-crystal

YUKAWA_THREADS=2 target/release/yukawa_crystal jel-controls crates/qf-yukawa-crystal
YUKAWA_THREADS=2 target/release/yukawa_crystal jel-main     crates/qf-yukawa-crystal 20
\end{lstlisting}
The initial-conditions script writes 20 periodic ($N=288$), 10 square-box ($N=289$) and 10 disk ($N=300$) CSV files
plus a \texttt{provenance.json}; every command except \texttt{refs} and \texttt{controls} reads them.
\texttt{main <dir> <kappa> <nrun>} runs \texttt{nrun} OC20-derived and \texttt{nrun} uniform starts, each quenched and
annealed-then-quenched with Hessian, plus \texttt{nrun/2} incommensurate-box controls; the counts of
Section~\ref{sec:num-results} correspond to $\texttt{nrun}=20$ at $\kappa=2$ ($2\cdot(20+20)+2\cdot10=100$) and
$\texttt{nrun}=10$ at $\kappa=1$ ($2\cdot(10+10)+2\cdot5=50$). \texttt{snapshot} re-runs selected published runs with
their seeds for the figures; \texttt{summarize.py} checks P1--P3 with the preregistered tolerances
(\texttt{REL\_\allowbreak TOL\_\allowbreak UNDERCUT = 1e-9}, \texttt{REL\_\allowbreak TOL\_\allowbreak LATTICE = 1e-8});
the jellium commands write under \texttt{results/jellium/}. Run times: the only times measured for this book are those of the unit tests (0.15~s for the ten tests, after
compilation) and of the Python companion (26.8~s). The production commands perform hundreds of CVODE relaxations of a
576-dimensional system with dense Jacobians and $3\times10^5$-step anneals; we did not time them here and quote no
figure. Start with \texttt{refs} and \texttt{controls}.

\section{The laboratory analysis script}\label{sec:num-lab}

The tiered paper~\cite{tiered} proposes laboratory tests with colloids and dusty plasmas (protocols L1--L4, \tierC{}
proposals; none has been run; Chapter~4 discusses the physics). They share
one analysis script, \texttt{papers/\allowbreak universal\_\allowbreak optimality\_\allowbreak tiered/\allowbreak lab/\allowbreak lab\_analysis.py}, which takes tracked 2D
positions and computes the same observables as the simulations: density from the mean Voronoi cell area of the
particles over which the energy is averaged; coordination with the $1.35\,d_0$ rule; local and global $|\psi_6|$; and the
interaction energy per interior particle for a stated law (\texttt{dipolar} $r^{-3}$, or \texttt{yukawa:<kappa>} with
$\kappa$ in the experiment's units), truncated at $r_c$ with a uniform tail correction
$\frac12\int_{r_c}^\infty2\pi r\,V(r)\,dr$ and compared with a triangular patch of the same density and truncation. Its
self-test feeds it four synthetic frames: a triangular patch in arbitrary units (scaled by $3.7$), the same with Gaussian
noise $0.08$, a square patch, and a Yukawa case that checks the unit conversion of $\kappa$. We ran
\texttt{taskset -c 2,3 python -I lab\_analysis.py --selftest}; exit code 0. The output is below; we left out fields,
marked \texttt{[...]} (our insertion), that repeat across blocks or are not discussed.
\begin{lstlisting}[language={},basicstyle=\ttfamily\scriptsize,title={lab\_analysis.py --selftest}]
{
 "triangular": {
  "n": 1600,
  "measured_density_input_units": 0.07304601899196492,
  "law": "dipolar",
  "law_density1": "dipolar",
  "rcut": 6.0,
  "counted": 728,
  "defect_fraction": 0.0,
  "psi6_local_mean": 1.0,
  "psi6_global": 1.0,
  "d0": 1.0745699318235395,
  [...]
  "e_per_particle": 4.464518999036898,
  "e_tri_same_density": 4.464518999036646,
  "rel_gap": 5.6399935141796596e-14,
  "prediction_e_ge_e_tri": true
 },
 "noisy_triangular": {
  "n": 1600,
  "measured_density_input_units": 1.0006342525142475,
  [...]
  "defect_fraction": 0.2625994694960212,
  "psi6_local_mean": 0.8244091272609662,
  "psi6_global": 0.8076197599391082,
  [...]
  "rel_gap": 0.029497434514365446,
  "prediction_e_ge_e_tri": true
 },
 "square": {
  "n": 1600,
  "measured_density_input_units": 1.0,
  [...]
  "defect_fraction": 1.0,
  "psi6_local_mean": 1.8369701987210297e-16,
  "psi6_global": 1.8369701987210538e-16,
  [...]
  "rel_gap": 0.011590527342694872,
  "prediction_e_ge_e_tri": true
 },
 "yukawa_units": {
  "n": 1600,
  "measured_density_input_units": 0.07304601899196492,
  "law": "yukawa:0.5405405405405405",
  "law_density1": "yukawa:2.0",
  [...]
  "defect_fraction": 0.0,
  [...]
  "e_per_particle_truncated": 0.39353696637141355,
  "e_tri_truncated": 0.3935369663713675,
  "tail_correction": 9.65130619565578e-06,
  [...]
  "rel_gap": 1.1707445433997635e-13,
  "prediction_e_ge_e_tri": true
 }
}
SELFTEST PASS
\end{lstlisting}
What the self-test shows, and what it does not:
\begin{itemize}
\item \emph{Unit invariance.} The triangular patch was given in units $3.7$ times too large (density $0.0730\ldots$ in
input units, about $1/3.7^2$); after rescaling, the defect fraction is $0$, $|\psi_6|=1$ and the relative gap is
$5.6\times10^{-14}$. The Yukawa case recovers $\kappa_1=2$ from $\kappa_{\rm input}=2/3.7$.
\item \emph{Thermal noise raises the energy.} The noisy patch has relative gap $+0.029$ and local $|\psi_6|=0.82$.
\item \emph{The defect detector fires on the square patch}: defect fraction $1.0$, $\psi_6\approx10^{-16}$.
\item \emph{What it does not show}: no self-test case returns \texttt{prediction\_e\_ge\_e\_tri = false}. Every input is a
configuration at or above the lattice energy, so the energy comparison is never shown to fire. A too-short cutoff
does make it fire (Exercise~\ref{ex:num-lab}).
\end{itemize}
The Yukawa value $0.393537$ (plain truncation at $r_c=6$) is not the same quantity as the shifted-force values $0.393231$
and $0.393533$ of Section~\ref{sec:num-python}: only equal truncation conventions may be compared.

\section{Exercises}\label{sec:num-ex}
All exercises are numerical experiments; their results are \tierX. Use at most two cores
(\texttt{taskset -c 2,3}). Before you report any number, run a positive and a negative control.

\begin{exercise}[Derivatives]\label{ex:num-grad}
Derive $\nabla E$ for the image-sum energy $E=\frac12\sum_{i,j}\sum_{s}{}'V_{\rm sf}(|x_i-x_j+s|)$ of the Python script
and show that the factor $\frac12$ disappears. Then check \texttt{energy\_grad} against central finite differences at
a random configuration (step $10^{-6}$), and verify Proposition~\ref{prop:grad-hess} by finite differences of the
gradient for one $2\times2$ block.
\end{exercise}

\begin{exercise}[Varying $\kappa$]
Run the Python script with $\kappa\in\{0.5,1,2,4\}$ at $n=6$, $r_c=8$. Record $e_{\rm tri}$ and how many of the five
random starts reach it. Is there a trend? What happens to the image-shell count $M$ and to the run time as
$\kappa\to0$ with $r_c$ fixed, and why must $r_c$ grow as $\kappa$ decreases for the truncation to stay at the same level?
\end{exercise}

\begin{exercise}[Varying $N$]
Run $n=4,5,6,7,8$ at $\kappa=2$. Plot the fraction of random starts that reach the lattice against $N=2n^2$, with more
seeds (say 20) for small $n$. Explain, in terms of the number of local minima, why the fraction should fall with $N$.
Is your result comparable with P3? List every difference between your experiment and the Rust study.
\end{exercise}

\begin{exercise}[Varying the cutoff]
At $n=6$, $\kappa=1$, compute $e_{\rm tri}$ with $r_c\in\{3,4,6,8,10\}$ and compare $|e_{\rm tri}(r_c)-e_{\rm tri}(10)|$
with $e^{-\kappa r_c}/r_c$ and with the unshifted sum of footnote~\ref{fn:unshifted}; explain why the shift dominates. Then replace the shifted-force potential by plain truncation ($V(r)$ for $r<r_c$, $0$
beyond) and observe what happens to L-BFGS: the energy is discontinuous at $r_c$. Why does that break a gradient-based
minimiser, and why does it not matter for the shifted-force form?
\end{exercise}

\begin{exercise}[A designed negative control]\label{ex:num-neg}
Reproduce the incommensurate-box negative control in Python. Put $N=n^2$ particles in a square box of side $n$
(density one) and relax from random starts. Check whether any end state is defect-free (at small $n$ a strained,
6-coordinated state is possible) and whether every one lies above the triangular value $e_{\rm tri}$ computed in a
commensurate box. Then plant a bug that affects only one side of a comparison, e.g.\ drop the factor $\frac12$ in
\texttt{energy\_grad} but not in \texttt{lattice\_sum}, and check that the box-sum versus direct-sum line of the output
now disagrees. (Dropping it everywhere would double $e$ and $e_{\rm tri}$ alike, and the comparison would not fire.) A
control that does not fire on a planted bug is not a control.
\end{exercise}

\begin{exercise}[Find the square-lattice saddle]\label{ex:num-square}
In the square box of side $n$ with $N=n^2$, place the perfect square lattice. (a) Show that $\max|\nabla E|$ is at
round-off level, so a gradient-based stopping rule accepts it. (b) Build the Hessian from \eqref{eq:Hblocks} and count
its negative eigenvalues with \texttt{numpy.linalg.eigvalsh} using the criterion $\lambda<-10^{-9}\lambda_{\max}$. The
Rust study found 72 at $\kappa=2$ and 64 at $\kappa=1$ for $N=289$, $r_c=8$; your count at smaller $n$ will differ.
(c) Perturb the lattice by $10^{-3}$ along the eigenvector of the most negative eigenvalue and relax. Where does it go,
and how far below $e_{\rm sq}$? (d) At the end state, check that exactly two eigenvalues are at round-off level, with
eigenvectors close to $\tau_x,\tau_y$; then put the particles in a harmonic bowl instead and explain why one zero mode
remains.
\end{exercise}

\begin{exercise}[L-BFGS versus gradient flow]
Integrate the gradient flow \eqref{eq:gflow} for the $n=4$ system with \texttt{scipy.integrate.solve\_ivp}, once with
\texttt{method="RK45"} and once with \texttt{method="BDF"} and the analytic Jacobian $-H$. Compare the number of
right-hand-side evaluations needed to reach $\max|F|<10^{-9}$. From the same random start, do L-BFGS and the gradient flow
end in the same state? Relate the RK45 cost to the stability bound $h<2/\lambda_{\max}$.
\end{exercise}

\begin{exercise}[Ewald by hand]\label{ex:num-ewald}
Write a Python version of \eqref{eq:ewald} for $N=72$ in the box of the script. With the minimum image you need
$r_c\le L_x/2\approx3.2$, so choose $\alpha$ large enough that $\operatorname{erfc}(\alpha r_c)$ is negligible, and
raise $k_{\rm cut}$ accordingly. (a) Check the Madelung energy of the
triangular lattice against $\eta\sqrt\pi$ with $\eta=-1.106102587$. (b) Check $\alpha$-invariance over a $4\times$ range,
scaling $k_{\rm cut}$ with $\alpha$. (c) Plant a bug: drop the background term $-N^2\sqrt\pi/(\alpha\mathcal A)$, and
show that the $\alpha$-invariance test fires.
\end{exercise}

\begin{exercise}[Make the lab self-test fire]\label{ex:num-lab}
Find an input for which \texttt{lab\_analysis.py} reports \texttt{prediction\_e\_ge\_e\_tri} as \texttt{false}, and
explain why this is a failure of the analysis assumptions and not of the theorem. Hint: a wrong $\kappa$ alone will not
do it, because the reference uses the same law. A cutoff shorter than the square lattice's second-neighbour distance
will: in a scratch call of \texttt{analyse} on the self-test's $40\times40$ square patch with the dipolar law and
$r_c=1.1$, we obtained a truncated energy $2.0$ against $2.4178$ for the triangular patch, and the flag \texttt{False}.
Why does the comparison invert?
\end{exercise}
